\paragraph{极小极大定理与零和博弈均衡}
设行玩家选 $i$、列玩家选 $j$ 时，列玩家收益为 $w_{i,j}$。双方的对偶线性规划为
\[\begin{array}{c@{\qquad}c}
\begin{aligned}
\max_{\mathbf p,v}\quad &v\\[-2pt]
\mathrm{s.t.}\quad &\sum_jw_{i,j}p_j\ge v\ (\forall i),\\[-2pt]
&p_j\ge0,\ \sum_jp_j=1
\end{aligned}
&
\begin{aligned}
\min_{\mathbf q,u}\quad &u\\[-2pt]
\mathrm{s.t.}\quad &\sum_iq_iw_{i,j}\le u\ (\forall j),\\[-2pt]
&q_i\ge0,\ \sum_iq_i=1.
\end{aligned}
\end{array}\]
极小极大定理保证 $u^*=v^*$；两边的最优策略共同构成纳什均衡。只解其中一边只能得到该玩家的最优策略，一般应使用线性规划而非贪心调整概率。
